\chapter{Conclusion générale}

\section{Synthèse du travail}

Ce mémoire a présenté \logminer{}, un prototype local destiné à analyser des journaux système et réseau hétérogènes à l'aide d'une architecture modulaire organisée autour d'agents fonctionnels. Le travail a porté sur l'ensemble de la chaîne de traitement : collecte, parsing, normalisation, routage vers des modèles adaptés, détection, priorisation d'anomalies candidates, corrélation, audit, mémoire adaptative et visualisation dans un tableau de bord.

La contribution principale n'est pas un nouvel algorithme d'apprentissage automatique. Elle se situe plutôt dans l'intégration contrôlée de plusieurs composants permettant de transformer des journaux de natures différentes en événements normalisés, en anomalies candidates et en incidents fonctionnellement exploitables par un analyste.

Le prototype distingue plusieurs familles de données, notamment Windows, Wazuh, Linux/auth, les flux réseau, HDFS, BGL et les entrées inconnues orientées vers une route de fallback. Cette organisation permet d'éviter l'application indistincte d'un modèle unique à des sources dont la structure et le sens opérationnel sont différents.

\section{Principaux résultats}

Les expérimentations montrent qu'une chaîne locale, légère et reproductible peut être mise en place pour traiter plusieurs familles de journaux. Les modèles supervisés obtiennent des résultats élevés dans certains protocoles exploratoires, mais les évaluations plus strictes rappellent que les performances dépendent fortement de la stratégie de séparation des données.

Le cas de CICIDS2017 est particulièrement important. Une séparation aléatoire produit un F1-score très élevé, tandis que les holdouts par fichier ou scénario font fortement diminuer le résultat. Cette observation montre qu'un score isolé ne suffit pas à démontrer une capacité de généralisation.

Les modèles non supervisés sont interprétés comme des générateurs d'anomalies candidates. Sur Wazuh, le système identifie 3 676 signaux parmi 122 563 événements. Ces signaux ne sont pas assimilés automatiquement à des intrusions confirmées ; ils constituent des éléments à filtrer, contextualiser et examiner.

Les campagnes Redis démontrent aussi l'intérêt architectural de l'approche par agents. Elles valident une répartition locale multiprocessus et la reprise de tâches non acquittées après panne simulée. La comparaison avec une chaîne monolithique montre cependant que cette architecture n'améliore pas automatiquement le débit sur de petits volumes. Son intérêt principal réside donc dans la modularité, l'observabilité, l'audit et la résilience locale.

Les mesures complémentaires de latence, de consommation CPU/RAM, de robustesse multiformat et de reproductibilité confirment que l'évaluation ne se limite pas aux scores prédictifs. Elles montrent que le prototype peut être exécuté, observé et comparé à partir d'artefacts techniques identifiables, même si ces mesures restent propres aux environnements et aux volumes testés.

\section{Réponse aux questions de recherche}

Les six questions de recherche formulées dans l'introduction peuvent être refermées de la manière suivante.

\begin{enumerate}

\item La normalisation des journaux hétérogènes est validée au niveau du prototype. Le schéma commun permet de rapprocher plusieurs sources tout en conservant les informations originales nécessaires à l'audit.

\item L'organisation en agents logiciels spécialisés est validée fonctionnellement. Elle améliore surtout la modularité, la traçabilité, l'observabilité et la reprise locale des tâches, sans constituer à elle seule une garantie de meilleur débit.

\item La sélection automatique du modèle est validée comme mécanisme architectural. Le routage par famille évite d'appliquer indistinctement un modèle unique à des données différentes, même si les expérimentations ne démontrent pas une supériorité prédictive systématique lorsque l'espace de caractéristiques reste commun.

\item Le tableau de bord répond fonctionnellement au besoin de consultation, de filtrage et de validation des anomalies candidates. En revanche, son utilisabilité auprès d'analystes réels n'a pas encore été évaluée par une étude utilisateur formelle.

\item La mémoire persistante est validée partiellement. Elle conserve les décisions, erreurs, durées et validations utiles à l'audit, à l'adaptation progressive du comportement des agents et à la mise à jour contrôlée des modèles. La priorisation a été implémentée sous la forme de mécanismes de sélection tenant compte de la famille des journaux, de l'historique des traitements et de certains retours enregistrés. Son efficacité propre n'a toutefois pas fait l'objet d'une évaluation quantitative indépendante.

\item Le protocole d'évaluation est validé comme cadre expérimental reproductible. Il combine des métriques de précision, rappel et F1, des mesures de latence, des observations CPU/RAM, des captures fonctionnelles, des scripts et des comparaisons contrôlées. Ses résultats doivent toutefois être interprétés dans les limites des jeux de données, des partitions et des environnements utilisés.

\end{enumerate}

\section{Mise à jour contrôlée des modèles}

Le système prévoit une mise à jour contrôlée des modèles en fonction des événements, de l'audit et des retours enregistrés par l'analyste. Les décisions de validation, de rejet ou de reclassement ne modifient pas directement un modèle en production. Elles enrichissent une mémoire exploitable lors des cycles suivants.

Le principe retenu est conservateur. Une nouvelle version d'un modèle est d'abord traitée comme un candidat. Elle est entraînée hors ligne, évaluée sur un jeu de validation déclaré, puis comparée au modèle courant avec une métrique définie à l'avance. Le candidat n'est promu que s'il obtient les meilleures performances dans le cadre retenu. L'ancien modèle reste conservé afin de permettre un retour arrière en cas de dégradation.

Cette organisation permet d'intégrer une logique d'amélioration progressive sans transformer le prototype en système d'auto-apprentissage non contrôlé. Elle maintient une séparation claire entre mémoire, décision humaine, entraînement candidat, validation et déploiement.

Dans ce cadre, le caractère autonome de \logminer{} doit être compris au sens défini dans l'introduction : le système peut orienter les journaux, sélectionner certains traitements, router les données vers les modèles disponibles et exécuter des tâches périodiques ou en micro-lots sans intervention manuelle constante. Cette autonomie reste encadrée par des règles explicites, par le registre de modèles et par la validation humaine.

De même, l'expression d'agents intelligents multitâches renvoie aux capacités déclarées, aux handlers spécialisés, à la mémoire locale, au heartbeat, à la politique de sélection des tâches et aux mécanismes de reprise. L'intelligence démontrée est donc une intelligence d'orchestration adaptative et vérifiable, et non une autonomie décisionnelle complète.

\section{Limites}

Plusieurs limites doivent être retenues.

Premièrement, les performances mesurées restent liées aux jeux de données, aux partitions et aux caractéristiques utilisées. Elles ne doivent pas être extrapolées directement à tous les environnements réels.

Deuxièmement, les expériences non supervisées produisent des anomalies candidates, mais pas une vérité terrain complète sur les incidents.

Troisièmement, la distribution validée est principalement locale et multiprocessus. Elle ne constitue pas encore une preuve de fonctionnement industriel multi-site.

Quatrièmement, le tableau de bord a été validé fonctionnellement, mais son utilisabilité auprès d'analystes réels reste à évaluer de manière plus structurée.

Enfin, la mise à jour contrôlée des modèles exige une discipline de suivi : conservation des métriques, séparation entre entraînement et validation, traçabilité des artefacts et surveillance de la dérive conceptuelle.

\section{Perspectives}

Les prolongements les plus importants concernent la validation multi-machine, l'amélioration de l'évaluation temporelle, l'anonymisation des journaux sensibles et l'étude de l'expérience utilisateur du tableau de bord.

Il serait également utile de renforcer les protocoles pour les journaux séquentiels comme HDFS et BGL, notamment par l'utilisation d'ensembles distincts d'entraînement, de validation et de test, ainsi que par une comparaison plus poussée avec des méthodes spécialisées dans les séquences de journaux.

Enfin, la mise à jour contrôlée des modèles pourra être enrichie par une meilleure structuration des retours analyste, une surveillance plus explicite de la dérive conceptuelle et une documentation systématique des promotions ou retours arrière de modèles.

\section{Conclusion finale}

\logminer{} montre qu'il est possible de construire une chaîne locale d'analyse de journaux à la fois modulaire, auditable et fonctionnellement exploitable par un analyste, sans que l'utilisabilité de l'interface ait encore été validée par une étude utilisateur formelle. Le prototype ne remplace pas un SOC industriel ni les plateformes de sécurité existantes, mais il fournit une base concrète pour comprendre, expérimenter et évaluer l'analyse de journaux hétérogènes.

L'intérêt du travail réside dans cette articulation entre architecture logicielle, modèles légers, agents multitâches, routage autonome encadré, mémoire adaptative, validation humaine et reproductibilité. Cette combinaison permet de traiter les anomalies comme des signaux à interpréter et à prioriser plutôt que comme des décisions automatiques définitives, ce qui correspond mieux aux exigences d'une application de cybersécurité.
